\section{Distinguishing a Moral Aspiration from a Causal Theory}

\subsection{The Hidden Mechanics of the Causal Chain}

A mission declares what contribution matters to the enterprise. It does not, by itself, explain how the organization will reliably generate that contribution.

Bridging this gap requires a \textbf{theory of change}: a falsifiable causal chain linking an operational intervention to an intended human outcome. While the phrase originated in non-profit and public-policy evaluation, the discipline is essential to every commercial business.

Consider the casual theory of change often recited by educational startups:

\begin{quote}
\itshape
If we select learners with adequate grit, diagnose their initial deficiencies, deliver structured conceptual explanations paired with deliberate practice, demand authentic capstone projects, and supply expert code reviews, then a higher proportion of them will develop the capacity to solve unfamiliar production engineering problems. If that capacity creates demonstrable economic value for employers or the learners themselves, customers will pay sustainable tuition for the training.
\end{quote}

Notice how much unexamined risk hides behind the word \textit{if}. Does the admissions screen genuinely predict persistence? Are the capstone rubrics pedagogically valid? Can expert instructors deliver rigorous diagnostic feedback without driving delivery costs above tuition revenue? Does the employer market recognize the resulting skill, or does it hire based on accredited university prestige?

A rigorous theory of change decomposes the optimistic arrow from ``great teaching'' to ``transformed careers'' into testable causal links, as mapped in Table~\ref{tab:theory-of-change}.

\begin{table}[htbp]
\centering
\begin{tcolorbox}[
  enhanced,
  colback=white,
  colframe=BayesBorder!80!white,
  boxrule=0.4pt,
  arc=2pt,
  left=3pt,right=3pt,top=3pt,bottom=3pt,
  before skip=3pt,after skip=4pt
]
\scriptsize
\setlength{\tabcolsep}{3pt}
\renewcommand{\arraystretch}{1.18}
\begin{tabularx}{\linewidth}{@{} >{\bfseries\color{BayesInk}}p{2.3cm} >{\color{BayesInk}}X @{}}
\toprule
\textbf{Causal Link \& Hypothesis} & \textbf{Observable Evidence vs Alternative Explanation} \\
\midrule
1. Inputs $\rightarrow$ Experience\newline
\textit{Instruction and exercises produce intended practice.} & 
\textbf{Evidence:} Cohort logs, code reviews, recorded sessions.\newline
\textbf{Rival Explanation:} Students learn independently from open documentation. \\
\midrule
2. Experience $\rightarrow$ Learning\newline
\textit{Deliberate practice and review build capability.} & 
\textbf{Evidence:} Blinded baseline and post-cohort diagnostic exams.\newline
\textbf{Rival Explanation:} Post-tests merely mimic practiced sample questions. \\
\midrule
3. Learning $\rightarrow$ Transfer\newline
\textit{Acquired skill functions in unfamiliar domains.} & 
\textbf{Evidence:} Complex workplace projects, novel technical tasks.\newline
\textbf{Rival Explanation:} High-ability students would have solved them anyway. \\
\midrule
4. Transfer $\rightarrow$ Value\newline
\textit{Employers or students capture tangible gain.} & 
\textbf{Evidence:} Promotions, raises, verified engineering productivity.\newline
\textbf{Rival Explanation:} Gains stem from macro hiring trends or preexisting networks. \\
\midrule
5. Value $\rightarrow$ Revenue\newline
\textit{Beneficiaries willingly fund the service.} & 
\textbf{Evidence:} Cash collections, organic referrals, corporate renewals.\newline
\textbf{Rival Explanation:} Early signups were driven by novelty or founder charisma. \\
\midrule
6. Revenue $\rightarrow$ Reinvestment\newline
\textit{Operating surplus expands institutional capacity.} & 
\textbf{Evidence:} Hiring research faculty, developing proprietary tools.\newline
\textbf{Rival Explanation:} Surplus is consumed by sales commissions and overhead. \\
\bottomrule
\end{tabularx}
\end{tcolorbox}
\caption{The 6 Causal Links in an Educational Theory of Change.}
\label{tab:theory-of-change}
\end{table}

This causal breakdown is indispensable even before empirical data exists. It shows the founder precisely which link in the chain requires testing. A mission that cannot be articulated as a series of plausible, testable links is an inspiring daydream, not an actionable foundation for an enterprise.

\subsection{The Reverse Incentive Diagnostic}

There is an equally vital diagnostic: identify what the organization's business model actually rewards leadership to optimize.

If an education business collects upfront tuition regardless of whether students complete the curriculum or master the concepts, the financial architecture contains a structural incentive to maximize marketing expenditure and minimize instructional overhead. This does not mean every education founder will surrender to cynical instincts. It does mean that if the mission is to survive, the founder must intentionally engineer countervailing measurement, admissions gating, pedagogical incentives, and governance checks.

\textbf{Every mission has an economic architecture, whether or not its founder wants to think about economics.}
